
\section{Empirical Strategy and Framework }
\label{section:strategy}

Oromia and Tigray fully implemented the 1994 language of instruction policy, and nearly all of their primary school students learn in their respective mother languages. However, in the Amhara state and The South, local languages are used as mediums of instruction only up to grades six and four, respectively, and all subjects are taught in English beyond these grades. In the capital city, Amharic is the language of instruction only up to grade six.\footnote{The four core states and the capital together account for about 90percent of Ethiopia's population, with the remaining peripheral states making up the rest.} 

The policy-induced disparity in MTI intake is a potential cause of differences in labor-market outcomes through its effects on education and mastery of the dominant language. The impact of MTI on employment outcomes can thus be identified by estimating equation (1), using the instrumental variables approach, where the exogenous instrument (Predicted MTI), the covariates, and the outcome variables are defined below. $MTIintake_{is} $ represents the number of years of schooling in the mother tongue of the student $ i $ in the state $ s $, and is influenced by several factors, including, crucially, the heterogeneous MTI laws in different states, which form the basis of the instrumental variable defined below. 
 
 %the interstate migration decisions of their families, and numerous other factors influencing decisions by students to drop out of schools before completing their primary education. 


\begin{equation}\label{eq:01}
	Y_{is} = \alpha + \beta\cdot MTIintake_{is} + X_{is}'\delta + \epsilon_{is}.
\end{equation}



The intake of MTI is endogenous in a regression of (1), as it can potentially be impacted by unobserved individual and family-level characteristics, which also affect the outcomes of the labor market. For example, although education is compulsory in Ethiopia for children between the ages of 5 and 16, attendance laws are barely enforced, and many students drop out before completing elementary school for many reasons, with serious consequences on their employment outcomes. 
 
Our identification strategy is also suitable for addressing issues that may arise due to noncompliance with assignment.  About 10percent of all students live outside their ethnic home states, and the majority of them do not have the option of learning in their ethnic languages. Although they constitute less than 1percent of the sample, there are also students who choose not to learn in their mother tongues and opt for the Amharic stream of education, likely expecting that education in Amharic might enhance their learning and economic outcomes later in life.  


The primary outcome variables, $ Y_{is} $, are employment outcomes for the student $ i $ in the state $ s $ in 2016, when the last complete YLS was conducted. They are indicator variables for salary employment and wage employment, which are switched on if subjects were employed when the survey was conducted in 2016.

The covariates $ X_{is}' $, include anthropomorphic measures of the students when they were eight years old, a measure of their stuntedness, an indicator of poverty based on the amount of food consumed, an index of family wealth, a measure of caregiver's education, family size, various indicators of asset ownership (including land, house, and consumer durables), access to electricity, an indicator of urbanization, a gender dummy, an indicator of type of school (public vs. private), and a few other relevant variables. 

\subsection*{The Instrumental Variable (IV)}

The instrumental variable, $ E_{is} $, defined in \eqref{eq:02}, is predicted MTI and is determined by the interaction of the ethnic backgrounds of the students and the MTI policies of the states where they are educated.

\begin{equation}\label{eq:02}
	E_{is} = \frac{N^{s}}{8}\cdot i^{M}.
\end{equation}

$ N^{s} $ is the number of legally mandated years of schooling in the mother tongue in the state $ s $; $ i^{M} $ is an indicator that is switched on if student $ i $ has access to mother tongue education. Some Amhara students who attend schools in certain school districts in Oromia and the South, and a small proportion of Oromo students who attend schools in the Amhara state, have the option to be educated in their native tongues. Consequently, in addition to students attending schools in their native tongues in their home states, $ i^{M} $ is switched on: 1) for Amhara students who learn in Amharic in some school districts in Oromia and The South, and 2) for Oromo students learning in Afaan Oromo in the special school districts of the Amhara state.

We argue that the IV is as good as randomly assigned and varies independently of job market outcomes. The only source of variation in the predicted MTI comes from the interaction of the ethnic backgrounds of the students and the educational policy of the state in which the students reside. 

The plausibility of the identification assumption that the predicted MTI affects the employment outcomes only through its effects on MTI intake can be tested in samples without obvious links between MTI intake and employment outcomes \citep{Angrist1990}. There is no persuasive reason why ethnicity would be the cause of differences in job market outcomes in the capital, since Amharic is the language of instruction in primary schools in the city, and almost all students are proficient in the language regardless of their ethnic backgrounds. Hence, we used the Addis Ababa sample for robustness checks.


%We also consider test scores (measures of cognitive ability) as intermediate variables through which MTI take-up is expected to influence labor market outcomes. We use two different test scores to capture variations in human capital accumulation by students schooled under the various regimes of education: the tests measure the students’ knowledge and aptitude in mathematics and verbal reasoning. 



























